\documentclass[11pt]{article}
\usepackage{amsmath,amssymb,amsthm}
\usepackage[margin=1in]{geometry}
\usepackage[T1]{fontenc}
\usepackage{hyperref}

\newtheorem{theorem}{Theorem}[section]
\newtheorem{lemma}[theorem]{Lemma}
\newtheorem{proposition}[theorem]{Proposition}
\newtheorem{corollary}[theorem]{Corollary}
\theoremstyle{remark}
\newtheorem{remark}[theorem]{Remark}

\title{Solution to Ramanujan Challenge Problem 2.1:\\
A Polynomial Continued Fraction for $6/(3-\pi)$}
\author{Xiang Huang\\{\small University of Illinois Springfield}\\{\small\texttt{xhuan5@uis.edu}}}
\date{July 2026}

\begin{document}
\maketitle

\begin{abstract}
We prove that the polynomial continued fraction of Problem~2.1 equals
$6/(3-\pi)$. The solution has two elementary ingredients and one classical
input. The elementary ingredients are (i) a polynomial index shift
$a_n = -\alpha(n+1)$, $b_n = \beta(n)$ identifying the challenge's coefficients
with those of Entry~5.3.22 of Cohen's dictionary of polynomial continued
fractions, and (ii) a sign-flip lemma: negating every partial denominator of a
continued fraction, while keeping the partial numerators, negates every
convergent. Because (ii) is proved \emph{at the level of convergents}, no
question of tail convergence arises. The classical input is the value of Cohen's
Entry~5.3.22 itself, namely $\pi = 3 + 6/T$ where $T$ is the corresponding
continued fraction; the challenge's fraction is exactly the sign-flip of $T$,
hence equals $-T = 6/(3-\pi)$. The argument is machine-checked in Lean~4, with
Entry~5.3.22 appearing as an explicit hypothesis.
\end{abstract}

\section{The problem}

Let
\begin{align}
a_n &= -220n^3 - 484n^2 - 301n - 42, \label{eq:an}\\
b_n &= 4n^2(2n+1)^2(5n-4)(5n+6). \label{eq:bn}
\end{align}
The challenge asks for a proof that
\begin{equation}\label{eq:main}
a_0 + \cfrac{b_1}{a_1 + \cfrac{b_2}{a_2 + \cfrac{b_3}{a_3 + \ddots}}}
= \frac{6}{3-\pi}.
\end{equation}

Throughout, for coefficient sequences $(c_n)_{n\ge0}$ and $(d_n)_{n\ge1}$ we
write $P_n,Q_n$ for the convergents of $c_0 + d_1/(c_1 + d_2/(c_2+\cdots))$,
defined by the classical recursions
\begin{equation}\label{eq:conv}
\begin{aligned}
P_{-1} &= 1, & P_0 &= c_0, & P_n &= c_nP_{n-1} + d_nP_{n-2},\\
Q_{-1} &= 0, & Q_0 &= 1,   & Q_n &= c_nQ_{n-1} + d_nQ_{n-2},
\end{aligned}
\end{equation}
and the value of the continued fraction means $\lim_n P_n/Q_n$ when it exists.

\section{The classical input}

\begin{theorem}[Cohen, Entry 5.3.22]\label{thm:cohen}
Let
\begin{equation}\label{eq:alphabeta}
\alpha(n) = 220n^3 - 176n^2 - 7n + 5,
\qquad
\beta(n) = 4n^2(2n+1)^2(5n-4)(5n+6).
\end{equation}
Then
\begin{equation}\label{eq:cohen}
\pi = 3 + \cfrac{6}{\alpha(1) + \cfrac{\beta(1)}{\alpha(2) +
\cfrac{\beta(2)}{\alpha(3) + \ddots}}}
= 3 + \cfrac{6}{42 + \cfrac{396}{1047 + \cfrac{38400}{4340 + \ddots}}}.
\end{equation}
\end{theorem}

This is Entry~5.3.22 of Cohen's \emph{Continued Fractions of Polynomial Type:
Theory and Encyclopedic Dictionary}~\cite{Cohen2026}, where it is recorded in
the machine-readable form
\begin{verbatim}
[()->Pi,[3,220*n^3-176*n^2-7*n+5],
         [6,4*n^2*(2*n+1)^2*(5*n-4)*(5*n+6)]]
\end{verbatim}
together with the displayed expansion whose first partial quotients are
$42, 396, 1047, 38400, 4340,\dots$; the same fraction appears as entry~1.2.18 of
the earlier database~\cite{Cohen2024}. Cohen obtains it from the hypergeometric
seed for $\pi = 6\arcsin(1/2)$ by a Bauer--Muir/Ap\'ery acceleration followed by
an even contraction and a polynomial equivalence transformation. This is the one
step of the present solution that we cite rather than reprove; everything else
below is proved here.

We have checked~\eqref{eq:cohen} numerically to $128$ decimal digits
(\S\ref{sec:numeric}), and verified the coefficients~\eqref{eq:alphabeta}
against the source entry, including its displayed values $42$, $396$, $1047$,
$38400$, $4340$.

\section{The index shift}

\begin{lemma}\label{lem:shift}
For all $n \ge 0$,
\[
a_n = -\alpha(n+1), \qquad b_n = \beta(n).
\]
\end{lemma}

\begin{proof}
The second identity holds because $b_n$ and $\beta(n)$ are the same expression.
For the first,
\begin{align*}
\alpha(n+1) &= 220(n+1)^3 - 176(n+1)^2 - 7(n+1) + 5\\
&= \bigl(220n^3 + 660n^2 + 660n + 220\bigr)
 - \bigl(176n^2 + 352n + 176\bigr) - 7n - 7 + 5\\
&= 220n^3 + 484n^2 + 301n + 42,
\end{align*}
so $-\alpha(n+1) = a_n$.
\end{proof}

Consequently, writing
\[
c_n := \alpha(n+1), \qquad d_n := \beta(n),
\]
the continued fraction $T := c_0 + d_1/(c_1 + d_2/(c_2+\cdots))$ is exactly the
tail appearing in~\eqref{eq:cohen}, and the challenge's fraction~\eqref{eq:main}
is the one with partial denominators $-c_n$ and the \emph{same} partial
numerators $d_n$.

\section{The sign-flip lemma}

\begin{lemma}\label{lem:signflip}
Let $(c_n),(d_n)$ be arbitrary, with convergents $P_n,Q_n$ as
in~\eqref{eq:conv}, and let $\widetilde P_n,\widetilde Q_n$ be the convergents
of the continued fraction with partial denominators $-c_n$ and partial
numerators $d_n$. Then for all $n \ge -1$,
\[
\widetilde P_n = (-1)^{n+1}P_n,
\qquad
\widetilde Q_n = (-1)^{n}Q_n .
\]
Consequently $\widetilde P_n/\widetilde Q_n = -\,P_n/Q_n$ for every $n$.
\end{lemma}

\begin{proof}
Induction on $n$. For the base cases,
$\widetilde P_{-1} = 1 = (-1)^0P_{-1}$,
$\widetilde P_0 = -c_0 = (-1)^1P_0$,
$\widetilde Q_{-1} = 0 = (-1)^{-1}Q_{-1}$ and
$\widetilde Q_0 = 1 = (-1)^0Q_0$.
For the step, assuming the claim at $n-1$ and $n-2$,
\[
\widetilde P_n = (-c_n)\widetilde P_{n-1} + d_n\widetilde P_{n-2}
= (-c_n)(-1)^{n}P_{n-1} + d_n(-1)^{n-1}P_{n-2}
= (-1)^{n+1}\bigl(c_nP_{n-1} + d_nP_{n-2}\bigr) = (-1)^{n+1}P_n,
\]
and identically
\[
\widetilde Q_n = (-c_n)(-1)^{n-1}Q_{n-1} + d_n(-1)^{n-2}Q_{n-2}
= (-1)^{n}\bigl(c_nQ_{n-1} + d_nQ_{n-2}\bigr) = (-1)^{n}Q_n .
\]
Finally
$\widetilde P_n/\widetilde Q_n = (-1)^{n+1}P_n/\bigl((-1)^nQ_n\bigr) = -P_n/Q_n$.
\end{proof}

We emphasise that Lemma~\ref{lem:signflip} is an identity for each fixed $n$.
It involves no hypothesis about convergence, and in particular no manipulation
of continued-fraction tails, which is where sign arguments of this kind usually
become delicate. It follows at once that the two continued fractions converge or
diverge together, and that when they converge their values are negatives.

\section{The main theorem}

\begin{theorem}\label{thm:main}
With $a_n,b_n$ as in~\eqref{eq:an}--\eqref{eq:bn},
\[
a_0 + \cfrac{b_1}{a_1 + \cfrac{b_2}{a_2 + \ddots}} = \frac{6}{3-\pi}.
\]
\end{theorem}

\begin{proof}
By Theorem~\ref{thm:cohen}, the continued fraction $T$ with partial denominators
$c_n = \alpha(n+1)$ and partial numerators $d_n = \beta(n)$ converges, and
$\pi = 3 + 6/T$; since $\pi \ne 3$ this gives
\[
T = \frac{6}{\pi - 3}.
\]
By Lemma~\ref{lem:shift}, the challenge's continued fraction has partial
denominators $a_n = -c_n$ and partial numerators $b_n = d_n$, so it is precisely
the sign-flip of $T$ in the sense of Lemma~\ref{lem:signflip}. By that lemma its
$n$-th convergent equals $-P_n/Q_n$ for every $n$, hence it converges, with
value
\[
-T = -\frac{6}{\pi-3} = \frac{6}{3-\pi}. \qedhere
\]
\end{proof}

\begin{remark}
Numerically $6/(3-\pi) = -42.3750798355862746187580309154\ldots$, and the
convergence is geometric at the golden-ratio rate: Cohen records convergence
type $\varphi^{-10}$ per term, i.e.\ $10\log_{10}\varphi = 2.0899\ldots$ decimal
digits per term, which is what we observe (\S\ref{sec:numeric}). In particular
the challenge's fraction converges at exactly the same rate as Cohen's, as
Lemma~\ref{lem:signflip} forces.
\end{remark}

\section{Formal verification}\label{sec:formal}

The accompanying Lean~4 development (\texttt{lean/}, toolchain
\texttt{leanprover/lean4:v4.30.0}) builds with \textbf{0 errors and
0 \texttt{sorry}} and formalises everything above except
Theorem~\ref{thm:cohen}:

\begin{center}
\begin{tabular}{ll}
\texttt{shift\_a}, \texttt{shift\_b} & Lemma~\ref{lem:shift}\\
\texttt{alphaC\_one}\,\dots\,\texttt{b21\_one} & the displayed values $42,396,1047,38400$\\
\texttt{cfP}, \texttt{cfQ} & the convergent recursions~\eqref{eq:conv}\\
\texttt{cfP\_neg}, \texttt{cfQ\_neg} & Lemma~\ref{lem:signflip}\\
\texttt{cf\_neg\_convergent} & $\widetilde P_n/\widetilde Q_n = -P_n/Q_n$, for every $n$\\
\texttt{challenge\_convergent\_eq} & the challenge is the sign-flip of Cohen's tail\\
\texttt{problem21\_pcf\_value} & Theorem~\ref{thm:main}, given Theorem~\ref{thm:cohen}\\
\end{tabular}
\end{center}

All of these depend only on the three standard axioms \texttt{propext},
\texttt{Classical.choice}, \texttt{Quot.sound}. There is no \texttt{native\_decide}.
Theorem~\ref{thm:cohen} enters as an explicit hypothesis:
\begin{verbatim}
theorem problem21_pcf_value
    (hCohen : Tendsto (fun k => cfP cohenC cohenD k / cfQ cohenC cohenD k)
                atTop (nhds (6 / (Real.pi - 3)))) :
    Tendsto (fun k => cfP a21 b21 k / cfQ a21 b21 k)
      atTop (nhds (6 / (3 - Real.pi)))
\end{verbatim}
so that the reliance on Cohen's entry is legible in the statement itself rather
than hidden in an axiom or a \texttt{sorry}.

\section{Numerical verification}\label{sec:numeric}

The script \texttt{verify.py} reimplements the same convergent
recursion~\eqref{eq:conv} used in the Lean file and checks, in exact rational
arithmetic where possible:

\begin{enumerate}
\item the index-shift identities of Lemma~\ref{lem:shift} for $0 \le n < 200$;
\item the coefficient values against Cohen's displayed expansion:
$\alpha(1)=42$, $\beta(1)=396$, $\alpha(2)=1047$, $\beta(2)=38400$,
$\alpha(3)=4340$, and $a_0=-42$, $a_1=-1047$, $b_1=396$;
\item the sign-flip identity $\widetilde P_n = (-1)^{n+1}P_n$,
$\widetilde Q_n = (-1)^{n}Q_n$ \emph{exactly}, for $0 \le n \le 41$;
\item convergence of~\eqref{eq:cohen} to $\pi$: the error at $n=61$ is
$5.4\times10^{-129}$;
\item convergence of the challenge's fraction to $6/(3-\pi)$: the error at
$n=61$ is $1.6\times10^{-126}$, matching the predicted $2.09$ digits per term.
\end{enumerate}

\begin{thebibliography}{10}

\bibitem{Cohen2024}
H.~Cohen,
\emph{A Database of Continued Fractions of Polynomial Type},
arXiv:2409.06086, 2024.

\bibitem{Cohen2026}
H.~Cohen,
\emph{Continued Fractions of Polynomial Type: Theory and Encyclopedic
Dictionary}, arXiv:2607.06581, 2026. Entry~5.3.22.

\bibitem{JonesThron1980}
W.~B.~Jones and W.~J.~Thron,
\emph{Continued Fractions: Analytic Theory and Applications},
Encyclopedia of Mathematics and its Applications~11, Addison--Wesley, 1980.

\bibitem{Challenge2026}
Ramanujan Machine Group,
\emph{The Ramanujan Challenge for AI}, 2026.

\end{thebibliography}

\end{document}
